% Compile with LuaLaTeX, XeLaTeX, or pdfLaTeX.
\documentclass[a4paper,11pt]{article}
\usepackage{iftex}
\ifPDFTeX
  \usepackage[T1]{fontenc}
  \usepackage[utf8]{inputenc}
  \usepackage{lmodern}
\else
  \usepackage{fontspec}
  \setmainfont{Latin Modern Roman}
\fi
\usepackage[margin=1.8cm]{geometry}
\usepackage{microtype,titlesec,enumitem,needspace,xcolor,fancyhdr,hyperref}
\usepackage{ragged2e}
\definecolor{linkcolor}{HTML}{164B73}
\definecolor{muted}{HTML}{505050}
\hypersetup{colorlinks=true,urlcolor=linkcolor,linkcolor=linkcolor,
  pdftitle={Kazuki Yano - CV},pdfauthor={Kazuki Yano}}
\urlstyle{same}
\setlength{\parindent}{0pt}
\setlength{\parskip}{0pt}
\raggedbottom
\pagestyle{fancy}
\fancyhf{}
\renewcommand{\headrulewidth}{0pt}
\fancyfoot[L]{\footnotesize\color{muted}Kazuki Yano}
\fancyfoot[R]{\footnotesize\color{muted}\thepage}
\setlength{\footskip}{20pt}
\titleformat{\section}{\large\bfseries}{}{0pt}{}[\vspace{1pt}\titlerule]
\titlespacing*{\section}{0pt}{10pt}{5pt}
\setlist[itemize]{leftmargin=1.15em,label=\textbullet,topsep=2pt,
  itemsep=0pt,parsep=0pt,partopsep=0pt}
% Dates use only the first row. Descriptions get the full text width.
% Each entry is kept together; line wrapping is left-aligned throughout.
\newcommand{\roleline}[3]{%
  \par\addvspace{5pt}\Needspace{4\baselineskip}%
  \noindent\textbf{#1}\hfill\mbox{\small #3}\par\nobreak
  {\itshape #2\par}\nobreak}
\newcommand{\publication}[4]{%
  \par\addvspace{5pt}%
  \begin{minipage}{\linewidth}\RaggedRight
    {\bfseries #1\par}
    {\small #2\par}
    {\small\textit{#3}\quad #4\par}
  \end{minipage}\par}
\newcommand{\award}[4]{%
  \par\addvspace{4pt}%
  \begin{minipage}{\linewidth}\RaggedRight
    \textbf{#1}\hfill\mbox{#2}\par
    {\small #3\par}
    {\small\itshape #4\par}
  \end{minipage}\par}
\begin{document}
\RaggedRight
% Avoid word fragments such as "Ex- / pansion" at line endings.
\hyphenpenalty=10000
\exhyphenpenalty=10000
\setlength{\emergencystretch}{1em}
\begin{center}
  {\fontsize{26}{30}\selectfont\bfseries Kazuki Yano}\par
  \vspace{5pt}
  {\small
  \href{mailto:yano.kazuki@dc.tohoku.ac.jp}{yano.kazuki@dc.tohoku.ac.jp}
  \quad\textbar\quad
  \href{https://kazuki-ya.github.io}{kazuki-ya.github.io}\par
  \vspace{2pt}
  \href{https://www.linkedin.com/in/kazuki-yano}{LinkedIn}
  \quad\textbar\quad
  \href{https://scholar.google.com/citations?user=KnttZUwAAAAJ&hl=en}{Google Scholar}}
\end{center}
\section{Summary}
I am a second-year Ph.D. student at Tohoku University working on large language model pre-training, training dynamics, and efficient model scaling, with recent research on multimodal foundation models and adapting text LLMs to new modalities.
I have first-author publications at leading venues including \textbf{ICLR 2026}, \textbf{INTERSPEECH 2026}, \textbf{COLM 2025}, and \textbf{NAACL 2025}.
I have hands-on experience in end-to-end large-scale model training, including pre-training, post-training, and controlled analysis on high-performance compute clusters.

\section{Education}
\textbf{Tohoku University}\hfill\mbox{Sendai, Japan}\par
\textbf{Ph.D. in Information Science} (expected)\hfill\mbox{2025 -- 2028}\par
\textbf{M.S. in Information Science}\hfill\mbox{2023 -- 2025}\par
{\small Supervisor for both degrees: \href{https://www.fai.cds.tohoku.ac.jp/members/js/}{Jun Suzuki}\par}
\roleline{Carnegie Mellon University}{Visiting Scholar}{Oct 2026 -- Mar 2027}
{\small WAVLab, Language Technologies Institute\par
Advisor: \href{https://sw005320.github.io/}{Shinji Watanabe}\par}

\section{Experience}
\roleline{Amazon Japan G.K.}{Research Fellow}{Jun 2026 -- Sep 2026}
\begin{itemize}
\item Developed LLM-based personalized recommendation, including user persona modeling and candidate reranking.
\end{itemize}
\roleline{SB Intuitions}{Research Intern}{Nov 2024 -- May 2026}
\begin{itemize}
\item Conducted LLM development on a large compute cluster, spanning pre-training and post-training.
\end{itemize}
\roleline{Preferred Networks Inc.}{Research Intern}{Aug 2024 -- Mar 2025}
\begin{itemize}
\item Built synthetic-data pipelines for LLM training; evaluated continued pre-training (CPT) and supervised fine-tuning (SFT).
\end{itemize}
\roleline{Rakuten Group, Inc.}{Data Scientist (Part-time)}{Jul 2024 -- May 2025}
\begin{itemize}
\item Studied representation/activation steering for style control via inference-time interventions.
\end{itemize}
\roleline{LINE Corporation}{Research Intern}{Aug 2023 -- Sep 2023}
\begin{itemize}
\item Developed an efficient adversarial-prompt optimization method for safety stress-testing aligned dialogue LMs.
\end{itemize}

\section{Publications}
% website-news: {"date": "2026-11", "location_en": "Athens, Greece", "location_ja": "アテネ🇬🇷"}
\publication{After the Predicate Comes the Phrase: How Inversion Affects LLM Reading Comprehension}
{Sumika Tsuda, \textbf{Kazuki Yano}, Jun Suzuki, and Reina Akama}
{Findings of EACL 2027}{\href{https://openreview.net/forum?id=wqGehzL9fy}{[Paper]}}
\publication{Adapting Text LLMs to Speech via Multimodal Depth Up-Scaling}
{\textbf{Kazuki Yano}, Jun Suzuki, and Shinji Watanabe}
{INTERSPEECH 2026}{\href{https://arxiv.org/abs/2604.00489}{[Paper]}}
\newpage
\section{Publications (continued)}
\publication{Pre-training LLM without Learning Rate Decay Enhances Supervised Fine-Tuning}
{\textbf{Kazuki Yano}, Shun Kiyono, Sosuke Kobayashi, Sho Takase, and Jun Suzuki}
{ICLR 2026}{\href{https://openreview.net/forum?id=JnebU2QLdH}{[Paper]}}
\publication{TimeMachine-bench: A Benchmark on Evaluating Model's Capability on Repository-level Migration Tasks}
{Ryo Fujii, Makoto Morishita, \textbf{Kazuki Yano}, and Jun Suzuki}
{EACL 2026}{\href{https://aclanthology.org/2026.eacl-long.385/}{[Paper]}}
\publication{Suppressing Final Layer Hidden State Jumps in Transformer Pretraining}
{Keigo Shibata, \textbf{Kazuki Yano}, Ryosuke Takahashi, Wataru Ikeda, and Jun Suzuki}
{Findings of EACL 2026}{\href{https://aclanthology.org/2026.findings-eacl.64/}{[Paper]}}
\publication{Efficient Construction of Model Family through Progressive Training \mbox{Using Model Expansion}}
{\textbf{Kazuki Yano}, Sho Takase, Sosuke Kobayashi, Shun Kiyono, and Jun Suzuki}
{COLM 2025}{\href{https://openreview.net/forum?id=fuBrcTH8NM}{[Paper]}}
\publication{Layer-wise Importance Analysis of Feed-Forward Networks in Transformer-based Language Models}
{Wataru Ikeda, \textbf{Kazuki Yano}, Ryosuke Takahashi, Jaesung Lee, Keigo Shibata, and Jun Suzuki}
{COLM 2025}{\href{https://openreview.net/forum?id=oP3b5YBFoP}{[Paper]}}
\publication{STEP: Staged Parameter-Efficient Pre-training for Large Language Models}
{\textbf{Kazuki Yano}, Takumi Ito, and Jun Suzuki}
{NAACL 2025}{\href{https://aclanthology.org/2025.naacl-short.32/}{[Paper]}}
\publication{Document-level Translation with LLM Reranking: Team-J at WMT 2024 General Translation Task}
{Keito Kudo, Hiroyuki Deguchi, Makoto Morishita, Ryo Fujii, Takumi Ito, Shintaro Ozaki, Koki Natsumi, Kai Sato, \textbf{Kazuki Yano}, Ryosuke Takahashi, Subaru Kimura, Tomomasa Hara, Yusuke Sakai, and Jun Suzuki}
{WMT 2024}{\href{https://aclanthology.org/2024.wmt-1.14/}{[Paper]}}
\publication{STEP: Staged Parameter-Efficient Pre-training for Large Language Models}
{\textbf{Kazuki Yano}, Takumi Ito, and Jun Suzuki}
{ACL 2024 Student Research Workshop}{}

\section{Honors and Awards}
\award{Outstanding Paper Award}{2025}
{The 31st Annual Meeting of the Association for Natural Language Processing (NLP 2025)}
{Awarded to top 14 / 765 submissions}
\award{BOOST Program for next-generation AI areas (JST)}{2025}
{Broadening Opportunities for Outstanding young researchers and doctoral students in STrategic areas}
{Research Fellow (subsidized by JST)}
\award{Encouragement Award}{2025}
{The 31st Annual Meeting of the Association for Natural Language Processing (NLP 2025)}
{Awarded to top 20 / 487 eligible submissions}
\award{Encouragement Award}{2024}
{The 30th Annual Meeting of the Association for Natural Language Processing (NLP 2024)}
{Awarded to top 18 / 427 eligible submissions}

\section{Academic Activities}
\textbf{Reviewer.}\quad ARR, TMLR, ICML (2026), NeurIPS (2026),
ICASSP (2026, 2027), ICLR (2027).\par
\smallskip
\textbf{Talk.}\quad
\href{https://www.wavlab.org/speech_lunch}
{Expanding Language Models for Speech and Beyond}.\par
{\small Speech Lunch, Carnegie Mellon University, Oct 8, 2026.\par}

\section{Skills}
\textbf{Tools.}\quad Python, PyTorch, distributed training, Linux, Git, Slurm, Docker/Singularity\par
\smallskip
\textbf{Languages.}\quad Japanese (native), English (professional working proficiency)
\end{document}
